\begin{table}[htbp]
\centering
\small
\caption{Urban Land Seizures in Santiago and Nationwide (1967--1971)}
\label{tab:urban_land_seizures_murphy}
\begin{tabularx}{\textwidth}{l c c c X}
\toprule
\textbf{Year} & \shortstack{\textbf{Santiago Land}\\\textbf{Seizures}} & \shortstack{\textbf{Chile Total}\\\textbf{Land Seizures}} & \shortstack{\textbf{Santiago}\\\textbf{Share (\%)}} & \textbf{Historical Phase} \\
\midrule
1967 & 13 & 8 & --- & Agrarian Reform \& Unionization Laws promulgated \\
1968 & 4 & 35 & 11.4\% & Consolidation of \textit{pobladores} committees \\
1969 & 23 & 220 & 10.5\% & Radicalization; Puerto Montt massacre \\
1970 & 103 & 560 & \textbf{18.4\%} & \textbf{Election of Salvador Allende (UP)} \\
1971 & 350 & --- & --- & \textbf{Reactivation boom \& urban mobilization} \\
\bottomrule
\end{tabularx}
\begin{minipage}{\textwidth}
\vspace{0.3em}
\footnotesize
\textit{Source:} \citet{Murphy2015}, p.~58. \textit{Note:} Demonstrates the exponential growth of the \textit{pobladores} movement and the spatio-temporal displacement of the class struggle into urban informality and collective social reproduction. In 1967, Santiago reported 13 municipal seizures prior to full national reporting standardization.
\end{minipage}
\end{table}
